% ---
% Agradecimentos
% ---
Agradeço à minha querida esposa \textbf{Amanda Ferrari}, minha maior inspiração para buscar o título de mestre. Sua presença constante, apoio incondicional e amor em todos os momentos — especialmente nos mais difíceis — foram fundamentais para que eu chegasse até aqui.

Agradeço à minha família: \textbf{minha mãe Ana Lúcia}, \textbf{meu padrasto Rosival Ronca}, \textbf{meu irmão Gustavo Ferrari}, \textbf{minhas avós Maria Lúcia e Offélia}, \textbf{minha tia Dalva Ávila}, \textbf{meu tio Luís Fernando Ávila}, \textbf{minha tia Heloísa Ferrari} e \textbf{meus primos Eloy Martins e Isadora Ávila}. Sem o apoio, incentivo e carinho de cada um de vocês, eu jamais teria tido a oportunidade de cursar uma universidade e trilhar este caminho acadêmico.

Ao meu querido pai \textbf{Marcos Sérgio Ferrari (in memoriam)} e à minha querida madrasta \textbf{Maria Cristina Studart (in memoriam)}, minha eterna gratidão e saudade.

Expresso também meu sincero agradecimento à minha orientadora, \textbf{Professora Doutora Maristela Oliveira dos Santos}, e à minha coorientadora, \textbf{Professora Doutora Eli Toso}. Este trabalho só foi possível graças à orientação, paciência e dedicação de ambas.

Agradeço, de forma especial, às minhas colegas de mestrado \textbf{Mestre Adriana Leticia dos Reis} e \textbf{Mestre Renata Cecconi de Souza}, pela disponibilidade, apoio e pela rica troca de experiências durante a jornada.

Agradeço, ainda, à minha amiga \textbf{Professora Doutora Christiane Lima Barbosa}, por seu apoio e incentivo nos primeiros passos do mestrado, e à \textbf{Professora Doutora Silmara Vicente}, com quem escrevi meu primeiro TCC e de quem recebi a semente da pesquisa científica, que floresceu e me trouxe até este momento.